\section{Fazit und Ausblick}
\label{sec:fazit}

% --------------------------------------------
% 6.1 Fazit
% --------------------------------------------
\subsection{Fazit}
\label{subsec:fazit}

Das Ziel dieser Arbeit war es, die Schritte der \ac{ip}-Coupon-Validierung zu untersuchen, redundante Zugriffe auf \textit{Railkit} zu identifizieren und diese durch einsetzen von Caching auf ein notwendiges Minimum zu reduzieren. Alles mit der Prämisse den verteilen Cache nicht unnötig zu belasten.

Dabei haben die Untersuchungen in Kapitel~\ref{sec:untersuchung} gezeigt, dass alle für die Coupon-Validierung benötigten Daten bereits nach der Angebotssuche vorliegen. Dennoch wurden drei \textit{Railkit} Anfragen zur Abfrage der Angebotsteile und zur Validierung der \textit{Offer} gestellt, um die bekannten Daten erneut zu laden.

Um diese Anfragen zu vermeiden wurde in Kapitel~\ref{sec:konzeption} eine Caching-Lösung entwickelt, die bereits bei der Angebotssuche den Cache mit Daten befüllt. Die Implementierung nach dem Cache-Aside-Muster speichert die validierungsrelevanten Werte von \textit{Railkit} in einem reduzierten Datenmodell ab. Ebenso bekommt jeder Eintrag eine mit \textit{Railkit} Angeboten abgestimmte Gültigkeit von zwei Stunden und Anfragen des Crawlers werden vom Caching ausgeschlossen, womit der verteilte Cache entlastet wird. Durch die Feststellung, dass die Mobile-App nicht das Prebooking durchläuft, fiel die Entscheidung den Cache bei der Angebotssuche statt beim Prebooking zu befüllen.

Die Evaluation in Kapitel~\ref{sec:evaluation} bestätigt, dass die Zielsetzung erreicht wurde:

\begin{itemize}
    \item \textbf{Externe Anfragen:} Die Aufrufe an \textit{Railkit} innerhalb der Coupon-Validierung entfallen vollständig. Somit  sinkt die Anzahl der Aufrufe bei CASA und Webshop von elf auf acht und in der Mobile-App von sechs auf drei über den Gesamten Buchungsprozess.
    \item \textbf{Antwortzeiten:} Die Antwortzeiten der Coupon-Operationen sinken im Median von 147 bis 634 Millisekunden auf zwei bis drei Millisekunden. Auch im 95. Perzentil liegen sie unter fünf Millisekunden.
    \item \textbf{Korrektheit:} Das Validierungsergebnis bleibt unverändert, da die Validierungsregeln nicht angepasst wurden und nur die Herkunft der Daten wechselt.
    \item \textbf{Aufwand:} Ein längere Verarbeitungszeit durch den zusätzlicher Schreibaufwand bei der Angebotssuche und Prebooking konnte nicht festgestellt werden. Ein Cache-Eintrag benötigt mit rund 1,1 KB nur etwa ein Zehntel des Umfangs der vollständigen Antwort von \textit{Railkit}.
\end{itemize}

Demhingegen steht als wesentlicher Kompromiss eine geringe Wiederverwendungsrate. Mit einem Verhältnis von rund 53 Angebotssuchen pro Buchung wird nur ein kleiner Teil der geschriebenen Einträge jemals gelesen. Aufgrund des reduzierten Datenmodells, der \ac{ttl} und dem Ausschluss des Crawlers, bleibt der Speicherbedarf jedoch begrenzt und liegt nach Abschätzungen im mittleren zweistelligen Megabyte-Bereich (vgl. Abschnitt~\ref{subsec:speicherbedarf}). Wobei dieser Wert aufgrund des einfachen Testfalls als untere Grenze zu verstehen ist. Am Ende überwiegt der Nutzen des effizienteren Validierungsprozesses dem geringen Aufwand pro Eintrag.

% --------------------------------------------
% 6.2 Ausblick
% --------------------------------------------
\subsection{Ausblick}
\label{subsec:ausblick}

Ausgehend von den Ergebnissen dieser Arbeit ergeben sich Ansatzpunkte für neue und weitereführende Untersuchungen und Verbesserungen am System der \ac{sbb}.

\paragraph{Überwachung im Produktivbetrieb}
Da die Evaluation nur über eine lokale Testumgebung mit einer einfachen Direkverbindung erfolgte, sollte die Wirksamkeit unter realen Bedinguen überwacht und bewertet werden. Mit einer protokollierung von Cache-Treffern und Fehlzugriffen ließe sich ermitteln und prüfen, wie häufig Validierung nach Ablauf der \ac{ttl} stattfinden. Ebenso sollte der Speicherverbrauch und die Verdrängung von Schlüsseln in Redis, zur Abschätzung des Speicherbedarfs, weiterhin überwacht werden.







\paragraph{Caching weiterer Prozesse}
Der Ansatz lässt sich auf weitere Prozesse übertragen, in denen bereits geladene Daten erneut von \textit{Railkit} abgefragt werden. Naheliegend ist das Prebooking: Hier werden das \textit{Offer} und die zugehörige Reiseroute erneut abgerufen, obwohl beide bereits aus der Angebotssuche vorliegen. Damit ließen sich zwei der fünf Aufrufe im Prebooking einsparen. Das reduzierte Datenmodell der Coupon-Validierung reicht dafür jedoch nicht aus, da das Prebooking deutlich mehr Informationen benötigt. Es wäre daher zu untersuchen, ob der Nutzen den höheren Speicherbedarf rechtfertigt, zumal das Prebooking im Gegensatz zur Coupon-Validierung bei jeder Buchung durchlaufen wird und die Wiederverwendungsrate damit höher ausfiele. Darüber hinaus könnten auch weitere Services des \ac{zvs} auf vergleichbare redundante Zugriffe hin untersucht werden.

\paragraph{Vergleich der Abläufe in Mobile-App, Webshop und CASA}
Die Untersuchung hat gezeigt, dass die Mobile-App den Coupon vor dem Prebooking validiert, während Webshop und CASA erst nach dem Prebooking validieren. Beide Varianten haben Vor- und Nachteile, die in einer weiteren Arbeit systematisch verglichen werden könnten:

\begin{itemize}
    \item \textbf{Validierung vor dem Prebooking (Mobile-App):} Der Kunde erhält früher eine Rückmeldung, ob sein Coupon gültig ist, und es entfallen die Abfrage der \textit{Order} sowie die Zuordnung von \textit{Order} zu \textit{Offers} im Cache. Allerdings ist zu diesem Zeitpunkt noch keine Kapazität reserviert, sodass sich Verfügbarkeit oder Preis bis zur Buchung noch ändern können.
    \item \textbf{Validierung nach dem Prebooking (Webshop, CASA):} Die Validierung erfolgt auf Grundlage einer bereits reservierten \textit{Order}, was den Stand der späteren Buchung genauer abbildet. Dafür ist ein zusätzlicher Prozessschritt notwendig, und die Validierung findet erst statt, nachdem bereits Kapazität bei \textit{Railkit} reserviert wurde.
\end{itemize}

Für die Caching-Lösung ist dieser Unterschied von besonderer Bedeutung. Würden alle Kanäle das Prebooking vor der Validierung durchlaufen, könnte der Cache wie in der ersten Iteration erst beim Prebooking befüllt werden. Bei einem Verhältnis von rund 53 Suchen pro Buchung würde dies die Anzahl der Cache-Einträge erheblich reduzieren. Würden umgekehrt alle Kanäle vor dem Prebooking validieren, könnte die Zuordnung von \textit{Order} zu \textit{Offers} entfallen. Eine Vereinheitlichung der Abläufe würde die Lösung somit vereinfachen, muss jedoch fachlich und aus Sicht der Nutzerführung bewertet werden.
